Multi-agent LLM systems exchange messages whose communicative necessity varies widely, yet no systematic evaluation protocol determines which messages can be safely removed. We introduce the first message-level functional analysis protocol for evaluating communication in two-agent LLM systems. The protocol comprises three reusable stages (annotate, perturb, prune), each deployable to new backbones, domains, and topologies without task-specific training data. We release a 21{,}661-message annotated corpus across 3~backbones and 4~reasoning domains with a five-class taxonomy for mechanism exploration, a zero-shot annotation pipeline using constrained decoding, and a statistical hypothesis-testing framework combining TOST equivalence, non-inferiority, and McNemar's tests. Human validation ($N{=}225$) confirms annotation reliability ($\kappa = 0.81$ binary ack/non-ack, underpinning all pruning claims; $0.76$ five-class); the cross-model annotation pipeline achieves $\kappa{=}0.72$. Applying the protocol across 18~configurations spanning 7B--72B scale, 4~reasoning domains, and 3~topology variants, with 10-seed replication for flagship settings, we surface a \emph{perturbation--pruning dissociation}: the most disruptive message type under single removal is safe to remove in bulk, enabling acknowledgment-targeted pruning that eliminates 36--56\% of inference calls at non-inferior accuracy under greedy decoding. The protocol further identifies boundary conditions where pruning fails and reveals backbone-dependent communication mechanisms, providing a screening tool for deployment decisions.
